\documentclass[11pt]{article}
\usepackage{amsmath,amssymb,amsthm}
\usepackage{bm}
\usepackage[margin=1in]{geometry}
\usepackage{hyperref}

\newtheorem{proposition}{Proposition}
\newtheorem{lemma}{Lemma}
\newtheorem{remark}{Remark}

\newcommand{\R}{\mathbb{R}}
\newcommand{\x}{\bm{x}}
\newcommand{\w}{\bm{w}}
\newcommand{\g}{\bm{g}}
\newcommand{\dbf}{\bm{\delta}}
\newcommand{\dW}{\Delta W}
\newcommand{\norm}[1]{\left\lVert #1 \right\rVert}
\newcommand{\ip}[2]{\langle #1, #2\rangle}
\newcommand{\sig}{\sigma}

\title{Why Linearization Accuracy and Adapter Leakage \emph{Decouple}:\\
a curvature ($\sig''$) account of a single control variable}
\author{Yoad Oxman}
\date{2026-08-13}

\begin{document}
\maketitle

\begin{abstract}
Empirically (STATUS.md, 2026-08-13) two quantities that ``ought'' to move together do not: making
the activation smoother \emph{lowers} the NTK/anchor linearization error (good for the
reconstruction method), yet on the LoRA/adapter path the \emph{kinked} ReLU --- the worst
linearizer --- \emph{leaks the most} instance-level information. This note explains why. The claim is
not that the two are accidentally uncorrelated, but that to leading order they are two readouts of the
\textbf{same} functional --- the activation's second derivative $\sig''$ evaluated at the
pre-activations --- with \emph{opposite} sign of benefit. Large $\sig''$ (a kink) makes the
per-sample gradient features \emph{diverse and separable} (leakage $\uparrow$) but the Taylor
expansion \emph{inaccurate} (lin-error $\uparrow$); small $\sig''$ (a smooth unit) does the reverse.
Hence no activation optimizes both; the attainable pairs lie on a frontier that the Softplus sharpness
$\beta$ dials directly. The account reproduces the observed effective-rank collapse
($\mathrm{rank}(\dW)\approx 1$ for Softplus), predicts why the effect is sharper on the LoRA path than
on full fine-tuning, and yields a cheap, falsifiable test (forward $+$ one backward pass, no
extraction).
\end{abstract}

\section{Setup}
Take the two-layer network used throughout the project,
\begin{equation}
  f(\x;\theta) \;=\; \sum_{j=1}^{m} v_j\,\sig\!\big(\ip{\w_j}{\x}\big),
  \qquad \theta = (W,\bm v),\quad W=[\w_1,\dots,\w_m]^\top\in\R^{m\times d}.
\end{equation}
A victim fine-tunes from a public $\theta_0$ on private data $\{(\x_i,y_i)\}_{i=1}^N$ by $T$ steps of
gradient descent on $L(\theta)=\sum_i \ell\!\big(f(\x_i;\theta),y_i\big)$ (BCE). The attacker observes
the weight change $\dW=\theta_T-\theta_0$ (full fine-tuning) or its low-rank LoRA factors $\dW=BA$,
and tries to recover $\{\x_i\}$.

\paragraph{Reconstruction.} Reconstruction linearizes $f$ around an anchor
$\theta_{\mathrm a}=(1-\alpha)\theta_0+\alpha\theta_T$ and matches $\dW$ to a weighted sum of
per-sample \emph{gradient features} $\phi(\x):=\nabla_\theta f(\x;\theta_{\mathrm a})$. To leading
order (exactly at $T{=}1$, and as the NTK term for $T>1$),
\begin{equation}\label{eq:obs}
  \dW \;\approx\; \sum_{i=1}^N c_i\,\phi(\x_i),
  \qquad c_i := -\eta\,\ell'\!\big(f(\x_i;\theta_{\mathrm a}),y_i\big)\in\R .
\end{equation}

\section{The observation is a low-rank factorization}
Restrict to the first layer (the reconstruction target). Since
$\partial f/\partial \w_j = v_j\,\sig'(\ip{\w_j}{\x})\,\x$, the first-layer gradient feature of a
sample is a \emph{rank-one} matrix,
\begin{equation}
  \phi_W(\x) \;=\; \g(\x)\,\x^\top,
  \qquad
  \g(\x) := \big(v_1\sig'(z_1),\dots,v_m\sig'(z_m)\big)^\top\in\R^m,\quad z_k:=\ip{\w_k}{\x}.
\end{equation}
The vector $\g(\x)$ is the \emph{gating pattern}: how strongly each neuron responds to $\x$.
Substituting into \eqref{eq:obs},
\begin{equation}\label{eq:factor}
  \boxed{\;\dW_1 \;=\; \sum_{i=1}^N c_i\,\g(\x_i)\,\x_i^\top \;=\; G\,\mathrm{diag}(\bm c)\,X^\top,\;}
  \qquad G=[\g(\x_1),\dots,\g(\x_N)]\in\R^{m\times N},\;\; X=[\x_1,\dots,\x_N]\in\R^{d\times N}.
\end{equation}
Reconstruction is exactly the problem of recovering the columns of $X$ from this rank-$\le\!N$
factorization.

\begin{proposition}[Identifiability $=$ conditioning of the gating matrix]\label{prop:ident}
The factorization \eqref{eq:factor} determines the individual $\x_i$ (up to the known per-column
scales $c_i$) \emph{iff $G$ has full column rank $N$}. If $\mathrm{rank}(G)=r<N$, only $r$ linear
combinations of the $c_i\x_i$ are observable; the remaining $N-r$ directions are a null space of
indistinguishable recombinations --- the \textbf{superposition problem}. Thus leakage is governed by
the conditioning $\kappa(G)=\sigma_1(G)/\sigma_N(G)$, \emph{not} by how accurate the linear model is.
\end{proposition}
\noindent (When $m<N$, $G$ cannot have rank $N$ and separation is impossible without extra structure;
this is the analogue for the low-rank LoRA case, \S\ref{sec:lora}.)

\section{The two quantities we care about}
\paragraph{(A) Linearization error} --- the project's headline anchor metric, the function-space
Taylor residual along the observed displacement $\dbf=\theta_T-\theta_{\mathrm a}$:
\begin{equation}
  \mathcal L_{\mathrm{lin}}
  \;:=\; \Big\| f(\x;\theta_{\mathrm a}+\dbf) - \big[f(\x;\theta_{\mathrm a}) + \phi(\x)^\top\dbf\big]\Big\|
  \;=\; \tfrac12\,\dbf^\top \nabla^2_\theta f(\x;\theta_\xi)\,\dbf + o(\norm{\dbf}^2).
\end{equation}
It measures how well the \emph{forward} map $\theta\mapsto f$ is approximated by its first-order model.

\paragraph{(B) Identifiability / leakage} --- by Proposition~\ref{prop:ident}, the conditioning
$\kappa(G)$ of the gating matrix; equivalently the effective rank of $\dW_1$. It measures whether the
\emph{inverse} problem ``$\dW\mapsto\{\x_i\}$'' is well posed.

\noindent These are different objects: (A) is forward-model accuracy, (B) is inverse-problem
conditioning. A model can approximate $f$ beautifully and still carry gradient features that are
nearly identical across samples (accurate but uninformative), or approximate $f$ terribly while
carrying sharply sample-specific features.

\section{The key lemma: both are first-order controlled by $\sig''$}
Everything hinges on one observation: the pre-activation curvature $\sig''$ simultaneously sets the
\emph{spread of the gates across samples} and the \emph{Taylor residual across the trajectory}.

\begin{lemma}[Gate diversity $\propto\sig''$]\label{lem:gate}
For two samples $\x_i,\x_j$ and neuron $k$, by the mean value theorem,
\begin{equation}
  \g_k(\x_i)-\g_k(\x_j) \;=\; v_k\big[\sig'(z_k^i)-\sig'(z_k^j)\big]
  \;=\; v_k\,\sig''(\zeta_k)\,\big(z_k^i-z_k^j\big),\qquad \zeta_k\in(z_k^j,z_k^i).
\end{equation}
Hence the column spread of $G$ --- and thus $\kappa(G)^{-1}$ and $\mathrm{rank}(G)$ --- scales to first
order with $\lvert\sig''\rvert$ at the pre-activations. If $\sig''\equiv0$ (affine unit) every column
of $G$ is identical: $\mathrm{rank}(G)=1$, total superposition. Large $\lvert\sig''\rvert$ spreads the
columns and lifts the rank toward $N$.
\end{lemma}

\begin{lemma}[Linearization residual $\propto\sig''$]\label{lem:lin}
The Hessian block is $\partial^2 f/\partial\w_k\partial\w_k^\top = v_k\,\sig''(z_k)\,\x\x^\top$, so
$\nabla^2_\theta f$ is linear in $\{\sig''(z_k)\}$ and
\begin{equation}
  \mathcal L_{\mathrm{lin}} \;=\; \tfrac12\sum_{k} v_k\,\sig''(\zeta_k)\,\big(\ip{\dbf_{\w_k}}{\x}\big)^2
  + (\text{cross terms}) \;=\; O\!\big(\lvert\sig''\rvert\cdot\norm{\dbf}^2\big).
\end{equation}
\end{lemma}

\begin{proposition}[The curvature tradeoff]\label{prop:tradeoff}
To leading order both the feature diversity that \emph{enables} reconstruction and the linearization
error that the NTK/anchor method must \emph{overcome} are proportional to the same functional, the
activation curvature $\sig''$ evaluated at the pre-activations:
\begin{equation}
  \text{leakage } \sim \kappa(G)^{-1} \;\uparrow\; \text{with } \lvert\sig''\rvert,
  \qquad\qquad
  \mathcal L_{\mathrm{lin}} \;\uparrow\; \text{with } \lvert\sig''\rvert .
\end{equation}
Consequently \textbf{no activation simultaneously maximizes leakage and minimizes linearization
error}; the attainable pairs lie on a frontier parameterized by $\sig''$. Smooth units
($\sig''\!\approx\!0$: GELU, SiLU, small-$\beta$ Softplus) occupy the \emph{low-error, low-leakage}
corner; the ReLU kink ($\sig''$ a Dirac mass at $0$, unbounded) occupies the \emph{high-error,
high-leakage} corner.
\end{proposition}

\noindent This is the whole story: the property that makes the Taylor expansion accurate --- $\sig'$
varying slowly --- is exactly the property that makes the gradient feature $\g(\x)=\sig'(\cdot)$
insensitive to \emph{which} $\x$ produced it. Smoothness buys an accurate model of a \emph{blurred}
signal.

\section{The Softplus-$\beta$ knob dials the frontier}
For $\sig_\beta(z)=\beta^{-1}\log(1+e^{\beta z})$ one has $\sig_\beta'(z)=s(\beta z)$ (logistic) and
\begin{equation}
  \sig_\beta''(z) \;=\; \beta\, s(\beta z)\big(1-s(\beta z)\big), \qquad \max_z \sig_\beta''=\tfrac\beta4 .
\end{equation}
So $\beta$ is a \emph{direct} dial on $\sig''$: $\beta\!\to\!0$ gives an affine unit ($\sig''\!\to\!0$,
maximal linearity, maximal superposition), $\beta\!\to\!\infty$ gives ReLU ($\sig''$ concentrates,
maximal identifiability, maximal lin-error). Proposition~\ref{prop:tradeoff} therefore makes a
\textbf{sharp one-parameter prediction}: along the existing \texttt{SOFTPLUS\_BETA\_SWEEP}, both
$\mathcal L_{\mathrm{lin}}$ and the leakage margin should increase monotonically with $\beta$, tracing
the frontier. This is the theory-cleanest experiment in the whole activation study.

\section{Consequences and consistency with the data}
\begin{itemize}
  \item \textbf{Effective-rank collapse (measured).} Proposition~\ref{prop:ident} predicts
  $\mathrm{rank}(\dW_1)\approx\mathrm{rank}(G)$: $\approx 1$ for a smooth unit, $\to N$ for ReLU. The
  rescore found Softplus \texttt{delta\_w\_effective\_rank}$=1$ vs.\ $2$ for the smooth family --- the
  smoothest reconstruction winner is exactly the most rank-collapsed, as predicted, and hence the
  least separable on the adapter path.
  \item \textbf{Why LoRA amplifies (\S\ref{sec:lora}).} The adapter records $P\dW$ for a rank-$r$
  projection $P$. Near-collinear (smooth) columns of $G$ collapse under $P$; distinct (ReLU) columns
  survive. So the identifiability gap between kinked and smooth \emph{widens} under the low-rank
  bottleneck --- which is why the empirical relu$>$smooth ordering is sharp on LoRA but muddier on
  full fine-tuning.
  \item \textbf{The anchor $\alpha$ is activation-dependent (observed).} Moving $\theta_{\mathrm a}$
  toward $\theta_T$ changes the pre-activations $z_k$ and hence $\sig''(z_k)$; for a smooth unit near
  saturation this can \emph{restore} some gate diversity (empirically the anchor rescues GELU/SiLU:
  LoRA margin $+0.03\!\to\!+0.14$), while it does nothing structural for a kink that is already
  maximally diverse (ReLU, best at $\alpha{=}0$) or for a curvature that is uniformly tiny.
\end{itemize}

\section{Falsifiable predictions and a cheap test}\label{sec:lora}
Both predictions need only a forward pass and one backward pass at $\theta_{\mathrm a}$ (\emph{no}
$50$k-epoch extraction):
\begin{enumerate}
  \item \textbf{Feature-rank vs.\ leakage.} Form $G=[\g(\x_1),\dots,\g(\x_N)]$ per activation and
  compute (i) the mean pairwise cosine of its columns and (ii) its numerical rank / $\kappa(G)$.
  Prediction: \emph{ReLU} $\Rightarrow$ low cosine, high rank (separable); \emph{smooth} $\Rightarrow$
  high cosine, rank $\approx 1$ (entangled). The rank should track the leakage margins while
  $\mathcal L_{\mathrm{lin}}$ tracks $-\!\sig''$.
  \item \textbf{Monotone $\beta$.} Along Softplus-$\beta$, both $\mathrm{rank}(G)$ and
  $\mathcal L_{\mathrm{lin}}$ increase with $\beta$.
\end{enumerate}
If (1)--(2) hold, the dissociation collapses to a one-line law: \emph{smoothness trades gradient-feature
identifiability for linearization accuracy, both governed by $\sig''$.} That is precisely the
project's open identifiability question (Q-A: when is the map $\dW\mapsto\{\x_i\}$ well posed?) with an
answer in terms of the model, not just the data.

\section{Caveats}
\begin{itemize}
  \item \emph{Homogeneity.} KKT/max-margin recovery (Haim et al.) assumes positively-homogeneous nets;
  smooth $\sig$ break homogeneity, so the relevant framework there is the NTK/gradient-feature account
  used here --- which needs no homogeneity. The factorization \eqref{eq:factor} is exact regardless.
  \item \emph{Leading order.} The ``$\propto\sig''$'' statements are first order for one hidden layer
  (MVT / second-order Taylor). Depth introduces products of gates, but the \emph{sign} of the effect
  (kink $\Rightarrow$ diverse features $+$ large residual) is preserved.
  \item \emph{Empirical status.} The supporting numbers are single-seed, $N\in\{2,10\}$, $T=10$,
  MNIST, oracle coefficients --- directional, not yet confirmed. The test in \S\ref{sec:lora} is what
  turns this note from hypothesis into measurement.
\end{itemize}

\end{document}
